\documentclass[eee-report]{otengineering}
% Later departmental changes can use table-start=0 (and figure/plate-start).
% \oteeereportsetup{left=1.5in,right=1in,table-start=0}
\title{Design and Implementation of a Monitoring System}
\author{John Doe \and Jane Doe}
\matricnumber{EEE/2023/001 and EEE/2023/002}
\projectdescription{A project report submitted to the Department of
 Electronic and Electrical Engineering in partial fulfilment of the
 requirements for the award of the Bachelor of Science degree}
\projectyear{2026}
\addbibresource{eee_report_references.bib}
\reportabbreviations{\begin{description}\item[ADC] Analogue-to-digital converter\end{description}}
\begin{document}
\maketitle
\frontmatter
\begin{acknowledgements}We thank our supervisors and colleagues.\end{acknowledgements}
\reportlists
\begin{abstract}This report demonstrates the departmental interface.\end{abstract}
\mainmatter
\chapter{Introduction}
\section{Background to the Study}
The design follows standard engineering practice~\cite{sample}.
\begin{quote}Long quotations retain the report's double spacing.\end{quote}
\begin{table}[htbp]\centering
\caption{Measured performance}\label{tab:performance}
\begin{tabular}{ll}\toprule Parameter & Value\\\midrule
Voltage & 5 V\\ Current & 250 mA\\\bottomrule\end{tabular}
\end{table}
\begin{figure}[htbp]\centering\fbox{System diagram}
\caption{System overview}\label{fig:overview}\end{figure}
\begin{plate}[htbp]\centering\fbox{Prototype photograph}
\caption{Completed prototype}\label{plate:prototype}\end{plate}
\chapter{Evaluation}
Chapter titles remain freely chosen. See Table~\ref{tab:performance}.
\begin{table}[htbp]\centering\caption{Second chapter measurement}
\begin{tabular}{ll}\toprule Test & Result\\\midrule A & Pass\\\bottomrule\end{tabular}\end{table}
\printbibliography
\appendix
\chapter{Supporting Measurements}
Appendices retain chapter-based numbering.
\end{document}
